%% sections/slide07-discussion.tex — SLIDE 7 — Discussion & limitations
%% =============================================================================
%% CONTENT BLUEPRINT — Discussion and Limitations
%% =============================================================================
%% Time budget: 60 seconds.
%%
%% WHAT TO PUT HERE:
%%
%%   LEFT half: Discussion (interpretation of the headline)
%%     • Two bullets connecting slide 5's number with theory:
%%         – "The measured value of $n_2 = 5.4 \pm 0.8 \times 10^{-19}$
%%            m²/W agrees with Sheik-Bahae's Z-scan relation (Eq. (3)) to
%%            within 12 %, supporting the assumption that ..."
%%         – "At higher intensities (≥3 GW/cm²), the deviation from
%%            lineart theory is consistent with two-photon absorption ..."
%%     • Be explicit about mechanism. Vague words ("because of thermal
%%       effects") do not help. Name the specific effect.
%%
%%   RIGHT half: Limitations (itemised, candid)
%%     • Maximum 4 limitations. Bullets, not prose.
%%     • Each bullet identifies:
%%         – What is limited (quantity / regime),
%%         – The magnitude of the limitation,
%%         – What it does NOT affect (proves the work still stands).
%%     • End the list with "None of these limitations affect the conclusion"
%%       if true. The examiner will appreciate the calm verdict.
%%
%% STYLE:
%%   • No new data on this slide. If you find yourself wanting to add a
%%     number, return to slide 5 or 6.
%%   • Limitations must be specific. "Sample purity" is vague; "Sample
%%     purity 99.9 %, supplier-stated, not independently verified" is
%%     specific.
%% =============================================================================

\begin{frame}{Discussion \& Limitations}
  \begin{columns}[T, totalwidth=\textwidth]
    \begin{column}{0.55\textwidth}
      \textbf{\color{sustBlueDark}Interpretation}\\[2pt]
      \begin{itemize}\setlength\itemsep{3pt}
        \item \textbf{The agreement is / is not consistent with [theory from
              slide 4]}. Specifically: \highlight{[One equation number or
              reference from slide 4]} predicts \highlight{[value]}, our
              measurement is \highlight{[value]}; $\Rightarrow$
              \highlight{[discrepancy percent]} discrepancy, explained by
              \highlight{[single most plausible mechanism]}.
        \item \textbf{The deviation at [extreme parameter] is consistent with}
              \highlight{[known second-order effect such as 2PA / thermal lens
              / exchange-correlation error / k-mesh sampling]}. The
              expected magnitude of this effect is \highlight{[X]}\,$\%$
              — matching our observation.
      \end{itemize}
      \vspace{4pt}
      \timenote{Cite slide 4 by equation number if you can; the examiner will
      cross-reference during Q\&A.}
    \end{column}

    \begin{column}{0.43\textwidth}
      \textbf{\color{sustBlueDark}Limitations}
      \begin{itemize}\setlength\itemsep{3pt}
        \item \textbf{[Limitation 1].} \itshape Specific magnitude. Does not
              affect Conclusion-1 (slide 8).
        \item \textbf{[Limitation 2].} \itshape Specific magnitude. Does not
              affect Conclusion-2.
        \item \textbf{[Limitation 3, optional].} \itshape Specific magnitude.
        \item \textbf{None of the above undermine headline result.}
      \end{itemize}

      \vspace{6pt}
      \begin{beamercolorbox}[sep=4pt]{block body example}
        \scriptsize\color{sustBlueDark}\itshape ``The principal uncertainty on
        the headline value of slide 5 is \highlight{X}\,\%, dominated by
        \highlight{[specific source named]}.''
      \end{beamercolorbox}
    \end{column}
  \end{columns}
\end{frame}
